% New English translation from checked cote 25, site 60 (PDF page 61).
% AI-translated by GPT-6 Astra (OpenAI, ChatGPT Pro mode).
\scanpage{25/60}
\var{\deleted{\illegible{2}}}{}{deletion}{ega25_25_060-v0001} \mspage{55}

To understand the meaning of this theorem, we should
\var{\inserted{\var{\deleted{\illegible{5}}}{}{deletion}{ega25_25_060-v0003}}}{\var{\deleted{\illegible{5}}}{}{deletion}{ega25_25_060-v0003}}{insertion}{ega25_25_060-v0002}
note from the beginning of the no. that if S=Spec(k), then for a
“sufficiently general” line L in $\hand{\check{\ }}P$, the regularity condition holds
(cf the catch-up additions mentioned in no. \hand{12}\var{\deleted{following}}{}{deletion}{ega25_25_060-v0005}\var{\inserted{\unsure{especially for S.3.|especially for 5.3.}}}{ \unsure{especially for S.3.|especially for 5.3.}}{insertion}{ega25_25_060-v0006}). Consequently, in the constructions
of “good” linear pencils indicated at the beginning of the no., we could
require the pencils in question \var{\xout{were}}{}{typewriter strike-out}{ega25_25_060-v0007} to satisfy that condition
(which \var{\xout{not}}{}{typewriter strike-out}{ega25_25_060-v0008} is a condition\var{\inserted[typed]{of the same type but}}{ of the same type but}{insertion}{ega25_25_060-v0009} distinct from requiring
that for every \hand{$\unsure{\xi|\zeta}\in$}L, $\hand{\unsure{\phi|\emptyset}}_{\hand{\unsure{\xi|\zeta}}}$ is $X_{k(\hand{\unsure{\xi|\zeta}})}$-regular). We should have included
the condition in question in the suggested summary statements.
\var{\xout{On the other hand, we may note}}{}{typewriter strike-out}{ega25_25_060-v0010} On the other hand, in practice 10.2 is
used chiefly in situation 9.5, which makes it desirable
to state the reformulation of 9.5 in terms of pencils only after
10.2, so that the statement in question can also include the
isomorphism of the pencil with a blow-up (i.e. give a description of the
situation that permits a convenient reference). We thus obtain a
way, for every smooth projective scheme \var{\inserted[typed]{geom \var{\xout{connected}}{}{typewriter strike-out}{ega25_25_060-v0013} conn X of di\unsure{m|p}$\hand{\geq}$2}}{geom \var{\xout{connected}}{}{typewriter strike-out}{ega25_25_060-v0013} conn X of di\unsure{m|p}$\hand{\geq}$2 }{insertion}{ega25_25_060-v0012}over an infinite field k, to find
a smooth closed subscheme\var{\xout{ everywhere}}{}{typewriter strike-out}{ega25_25_060-v0016} \var{\inserted[typed]{nonempty}}{nonempty}{insertion}{ega25_25_060-v0017} of codimension \unsure{\var{\deleted{\illegible{1}}}{}{deletion}{ega25_25_060-v0018}|\hand{$\geq$}}2 at each of
\var{\xout{can be}}{}{typewriter strike-out}{ega25_25_060-v0019}
its points \var{\xout{(which does not exclude its being empty) and it implies on the}}{}{typewriter strike-out}{ega25_25_060-v0020}
\var{\xout{contrary if dim X}}{}{typewriter strike-out}{ega25_25_060-v0021} such that the blow-up scheme admits a fibration
over P$^1$, all of whose\hand{/}fibers \var{\inserted[typed]{are geom integral, and all of whose fibers}}{are geom integral, and all of whose fibers }{insertion}{ega25_25_060-v0022}are smooth are at most a finite number,
the latter having at most one \var{\xout{nonsmooth (and}}{}{typewriter strike-out}{ega25_25_060-v0023} geom. singular point, and
such a point \var{\xout{is}}{}{typewriter strike-out}{ega25_25_060-v0024} \var{\inserted[typed]{being}}{being}{insertion}{ega25_25_060-v0025} rational over k and geom. ordinary singular. This
explains the importance of a thorough study (\var{\deleted{barely}}{}{deletion}{ega25_25_060-v0026} \var{\inserted{just}}{just}{insertion}{ega25_25_060-v0027} begun at
present) of such fibrations with singular fibers, to reduce to
some extent the study of smooth projective varieties of dim d
to that of (one-parameter families of) projective varieties of dim d-1\hand{,}
\hand{which may have ordinary singularities.}
\tnote{25-060-1}{The French stray-letter and agreement cancellations in rattrappages, description, trouver and un give no independent English corrections. The source literally has sont lisses sont au plus un nombre fini; the repeated are is retained without replacing either word by an unstated except. The uncertain codimension alternative and the abbreviated conn remain as written.}

\var{\margin{60}}{\placedmargin{60}}{marginal note placed}{ega25_25_060-v0029}
